% Options for packages loaded elsewhere
\PassOptionsToPackage{unicode}{hyperref}
\PassOptionsToPackage{hyphens}{url}
\PassOptionsToPackage{dvipsnames,svgnames,x11names}{xcolor}
\documentclass[
  spanish,
  10pt,
  a4paper,
]{article}
\usepackage{xcolor}
\usepackage[top=2.2cm,bottom=2.2cm,left=2.1cm,right=2.1cm]{geometry}
\usepackage{amsmath,amssymb}
\setcounter{secnumdepth}{5}
\usepackage{iftex}
\ifPDFTeX
  \usepackage[T1]{fontenc}
  \usepackage[utf8]{inputenc}
  \usepackage{textcomp} % provide euro and other symbols
\else % if luatex or xetex
  \usepackage{unicode-math} % this also loads fontspec
  \defaultfontfeatures{Scale=MatchLowercase}
  \defaultfontfeatures[\rmfamily]{Ligatures=TeX,Scale=1}
\fi
\usepackage{lmodern}
\ifPDFTeX\else
  % xetex/luatex font selection
\fi
% Use upquote if available, for straight quotes in verbatim environments
\IfFileExists{upquote.sty}{\usepackage{upquote}}{}
\IfFileExists{microtype.sty}{% use microtype if available
  \usepackage[]{microtype}
  \UseMicrotypeSet[protrusion]{basicmath} % disable protrusion for tt fonts
}{}
\makeatletter
\@ifundefined{KOMAClassName}{% if non-KOMA class
  \IfFileExists{parskip.sty}{%
    \usepackage{parskip}
  }{% else
    \setlength{\parindent}{0pt}
    \setlength{\parskip}{6pt plus 2pt minus 1pt}}
}{% if KOMA class
  \KOMAoptions{parskip=half}}
\makeatother
\usepackage{longtable,booktabs,array}
\newcounter{none} % for unnumbered tables
\usepackage{calc} % for calculating minipage widths
% Correct order of tables after \paragraph or \subparagraph
\usepackage{etoolbox}
\makeatletter
\patchcmd\longtable{\par}{\if@noskipsec\mbox{}\fi\par}{}{}
\makeatother
% Allow footnotes in longtable head/foot
\IfFileExists{footnotehyper.sty}{\usepackage{footnotehyper}}{\usepackage{footnote}}
\makesavenoteenv{longtable}
\ifLuaTeX
\usepackage[bidi=basic,shorthands=off]{babel}
\else
\usepackage[bidi=default,shorthands=off]{babel}
\fi
\ifLuaTeX
  \usepackage{selnolig} % disable illegal ligatures
\fi
\setlength{\emergencystretch}{3em} % prevent overfull lines
\providecommand{\tightlist}{%
  \setlength{\itemsep}{0pt}\setlength{\parskip}{0pt}}
% Shared visual layer for the generated Kaleido FlowTwin reports.
\usepackage{helvet}
\renewcommand{\familydefault}{\sfdefault}
\usepackage{microtype}
\usepackage{xcolor}
\usepackage{titlesec}
\usepackage{fancyhdr}
\usepackage{enumitem}
\usepackage{booktabs}
\usepackage{xurl}
\usepackage{lastpage}
\usepackage{etoolbox}

\definecolor{KaleidoNavy}{HTML}{0A2342}
\definecolor{KaleidoBlue}{HTML}{176B87}
\definecolor{KaleidoCyan}{HTML}{3BB7C4}
\definecolor{KaleidoMint}{HTML}{DDF5F1}
\definecolor{KaleidoFog}{HTML}{F2F5F7}
\definecolor{KaleidoSlate}{HTML}{5D6A7A}
\definecolor{KaleidoCoral}{HTML}{F26B5B}

\titleformat{\section}
  {\Large\bfseries\color{KaleidoNavy}}
  {\thesection}{0.75em}{}[\vspace{0.2em}\color{KaleidoCyan}\titlerule]
\titleformat{\subsection}
  {\large\bfseries\color{KaleidoBlue}}
  {\thesubsection}{0.65em}{}
\titleformat{\subsubsection}
  {\normalsize\bfseries\color{KaleidoNavy}}
  {\thesubsubsection}{0.55em}{}
\titlespacing*{\section}{0pt}{2.2em}{1em}
\titlespacing*{\subsection}{0pt}{1.6em}{0.65em}

\pagestyle{fancy}
\fancyhf{}
\fancyhead[L]{\small\bfseries\color{KaleidoNavy}KALEIDO FLOWTWIN}
\fancyhead[R]{\small\color{KaleidoSlate}DOSSIER TECNICO}
\fancyfoot[L]{\scriptsize\color{KaleidoSlate}EVOCON Solutions GmbH}
\fancyfoot[R]{\scriptsize\color{KaleidoSlate}\thepage/\pageref*{LastPage}}
\renewcommand{\headrulewidth}{0.6pt}
\renewcommand{\headrule}{\hbox to\headwidth{\color{KaleidoCyan}\leaders\hrule height \headrulewidth\hfill}}
\renewcommand{\footrulewidth}{0pt}
\setlength{\headheight}{20pt}

\setlist[itemize]{leftmargin=1.4em,itemsep=0.25em,topsep=0.35em}
\setlist[enumerate]{leftmargin=1.7em,itemsep=0.25em,topsep=0.35em}
\setlength{\parindent}{0pt}
\setlength{\parskip}{0.55em}
\setlength{\emergencystretch}{3em}
\renewcommand{\arraystretch}{1.18}

\AtBeginEnvironment{quote}{%
  \color{KaleidoNavy}\itshape
  \setlength{\leftskip}{1em}
  \setlength{\rightskip}{1em}
}

\makeatletter
\renewcommand{\maketitle}{%
  \hypersetup{pageanchor=false}
  \begin{titlepage}
    \pagecolor{KaleidoNavy}
    \color{white}
    \vspace*{1.4cm}
    {\color{KaleidoCyan}\rule{2.4cm}{2.5pt}\par}
    \vspace{1.1cm}
    {\Huge\bfseries\@title\par}
    \vspace{0.65cm}
    {\Large\color{KaleidoMint}Kaleido FlowTwin\par}
    \vfill
    {\large Álvaro Schwiedop Souto\par}
    \vspace{0.15cm}
    {\normalsize\color{KaleidoMint}Industrial AI \& Predictive Quality Lead\par}
    \vspace{0.15cm}
    {\normalsize\color{KaleidoMint}EVOCON Solutions GmbH\par}
    \vspace{0.25cm}
    {\normalsize\color{KaleidoMint}\@date\par}
  \end{titlepage}
  \hypersetup{pageanchor=true}
  \nopagecolor
  \color{black}
}
\makeatother
\usepackage{bookmark}
\IfFileExists{xurl.sty}{\usepackage{xurl}}{} % add URL line breaks if available
\urlstyle{same}
% fallback for those not using the hyperref driver hyperxmp:
\makeatletter
\@ifundefined{xmpquote}{\newcommand{\xmpquote}[1]{#1}}{}
\makeatother
\hypersetup{
  pdftitle={Auditoria de reutilizacion de proyectos},
  pdflang={es},
  colorlinks=true,
  linkcolor={KaleidoBlue},
  filecolor={Maroon},
  citecolor={Blue},
  urlcolor={KaleidoBlue},
  pdfcreator={LaTeX via pandoc}}

\title{Auditoria de reutilizacion de proyectos}
\author{}
\date{15 julio 2026}

\begin{document}
\maketitle

{
\setcounter{tocdepth}{3}
\tableofcontents
}
\section{Auditoria de reutilizacion de
proyectos}\label{auditoria-de-reutilizacion-de-proyectos}

Fecha: 15-07-2026. La evaluacion distingue codigo implementado,
especificacion y resultado promocionable.

\subsection{Veredicto}\label{veredicto}

No existe un proyecto que pueda presentarse a Kaleido como MVP ya
ajustado. La mejor estrategia es crear un repositorio nuevo y reutilizar
componentes auditados, no copiar un proyecto completo.

{\def\LTcaptype{none} % do not increment counter
\begin{longtable}[]{@{}llll@{}}
\toprule\noalign{}
Proyecto & Encaje & Reutilizar & No reutilizar como evidencia \\
\midrule\noalign{}
\endhead
\bottomrule\noalign{}
\endlastfoot
\texttt{predictiveops-worldmodel} & Alto como esqueleto & contratos
action/context/outcome, leakage, splits, lead time, SOTA gate, tests &
la narrativa CNC universal ni el modelo actual como predictor
logistico \\
\texttt{industrial\_jepa\_mvp} & Medio como biblioteca de ideas &
agregacion jerarquica, reportes, baselines, patrones de dashboard &
todas las metricas pre-audit; actualmente 0 resultados claim-eligible \\
\texttt{e-jepa-ttc} & Bajo/medio & streaming, incertidumbre,
multi-horizonte, cache y export futuro & modelo de camara de eventos;
test reutilizado y sin resultado post-fix \\
\texttt{softnav-jepa} & Bajo para el piloto & separacion
planner/constraints, candidate ranking, safety fallback & resultados
v1-v3 invalidados; proxy cinematico, no puerto/robot real \\
\texttt{SemanticSegmentation3Dclouds} & Futuro 3D & schemas geograficos,
split espacial, metricas, LiDAR & DINO/Point-JEPA como mejora probada;
trabajo sucio y bloques corregidos pendientes \\
\texttt{LiDAR\ Foundational\ Model} & Futuro 3D & plan ALS-first y gates
de financiacion & codigo/modelo: hoy es especificacion, no
implementacion \\
\texttt{LeJEPAenAQilles} & Nulo para producto & diagnosticos de colapso
y disciplina experimental & smoke colapsado y dominio molecular \\
\texttt{HAS-JEPA} & Ninguno & nada & es una nota de transcripcion
musical, no un proyecto logistico \\
\end{longtable}
}

\subsection{\texorpdfstring{\texttt{predictiveops-worldmodel}}{predictiveops-worldmodel}}\label{predictiveops-worldmodel}

Estado verificado:

\begin{itemize}
\tightlist
\item
  88 ficheros de proyecto;
\item
  43 tests pasan;
\item
  contratos de columnas y validacion action/context implementados;
\item
  adaptadores, leakage, splits, timestamps, lead time y SOTA gate
  implementados;
\item
  modelo Sensor-JEPA experimental sobre CNC;
\item
  no hay conectores logisticos, OCEL/XES, Trace Port ni simulacion
  discreta;
\item
  las configuraciones industriales avanzadas descritas en README son en
  gran parte roadmap, no implementacion.
\end{itemize}

Decision: usarlo como referencia de gobernanza. Portar o extraer modulos
pequenos con tests y licencia, manteniendo un historial de procedencia.
No hacer fork ciego porque su semantica y datasets son industriales/CNC.

\subsection{\texorpdfstring{\texttt{industrial\_jepa\_mvp}}{industrial\_jepa\_mvp}}\label{industrial_jepa_mvp}

Estado verificado:

\begin{itemize}
\tightlist
\item
  demos de sensores CNC y anomalia visual;
\item
  ramas de DenseSensorJEPA, token world model, DINO/PatchCore/PaDiM y
  jerarquia;
\item
  historial de fugas: hardness repetida entre splits y seleccion de
  umbral con test;
\item
  \texttt{STATUS.md} declara 0 resultados post-audit elegibles;
\item
  worktree con artefactos no versionados.
\end{itemize}

Decision: reutilizar ideas y, tras revision de licencia/API, funciones
de agregacion/reporting. Nunca copiar cifras a la presentacion.

\subsection{\texorpdfstring{\texttt{e-jepa-ttc}}{e-jepa-ttc}}\label{e-jepa-ttc}

Estado verificado:

\begin{itemize}
\tightlist
\item
  pipeline real de camara de eventos, representaciones, temporal JEPA y
  TTC;
\item
  nueve secuencias locales, pero CPLA-high ya fue inspeccionado;
\item
  \texttt{STATUS.md} declara 0 metricas post-fix promovibles;
\item
  faltan robustness, ONNX, streaming demo final y report builder
  completo.
\end{itemize}

Decision: tomar patrones de ventanas temporales, multi-horizonte e
incertidumbre. El dominio visual asincrono no se traslada a eventos de
negocio.

\subsection{\texorpdfstring{\texttt{softnav-jepa}}{softnav-jepa}}\label{softnav-jepa}

Estado verificado:

\begin{itemize}
\tightlist
\item
  buena arquitectura hibrida: stack clasico, restricciones duras, JEPA
  local;
\item
  datasets v1-v3 y resultados derivados invalidados por inputs
  privilegiados, split/policy y perturbaciones no emparejadas;
\item
  v4 corrige contratos, pero no se ha ejecutado el benchmark completo;
\item
  no hay evidencia de Isaac, ROS, MuJoCo ni robot real.
\end{itemize}

Decision: conservar el principio
\texttt{learned\ ranking\ never\ overrides\ hard\ constraints}. Solo
aplicaria en una futura fase de automatizacion fisica.

\subsection{\texorpdfstring{\texttt{SemanticSegmentation3Dclouds} y
\texttt{LiDAR\ Foundational\ Model}}{SemanticSegmentation3Dclouds y LiDAR Foundational Model}}\label{semanticsegmentation3dclouds-y-lidar-foundational-model}

Estado verificado:

\begin{itemize}
\tightlist
\item
  corpus ALS grande y split geografico disciplinado;
\item
  baseline geometrico fuerte; DINO tardio no aporta mejora robusta;
\item
  Point-JEPA actual usa centros por muestreo uniforme aproximado, no FPS
  real;
\item
  estado declara ausencia de bloques procesados corregidos y pesos DINO
  reales;
\item
  GeoLiDAR-FM es un plan ALS-first, aun sin implementacion.
\end{itemize}

Decision: linea futura para inventario/ocupacion/seguridad espacial del
patio o terminal, solo si Kaleido formula un problema 3D y
aporta/sensoriza datos.

\subsection{Componentes que debe crear
FlowTwin}\label{componentes-que-debe-crear-flowtwin}

\begin{itemize}
\tightlist
\item
  adaptador Trace Port/API/CSV;
\item
  esquema OCEL-like para proyecto, operacion, turno, carga y recurso;
\item
  versionado de plan y replanificacion;
\item
  process discovery/conformance;
\item
  predictor de tiempo restante y riesgo calibrado;
\item
  simulador discreto sencillo;
\item
  capa de escenarios con acciones verificadas;
\item
  dashboard para prefijos vivos y explicacion;
\item
  evaluacion comercial y tecnica por proyecto/turno.
\end{itemize}

\subsection{Regla de procedencia}\label{regla-de-procedencia}

Cada modulo copiado debe registrar repositorio, commit, licencia,
fichero original, cambios y tests. Ningun resultado historico se copia:
todos los resultados de FlowTwin se regeneran desde datos FlowTwin y
artefactos propios.

\end{document}
